% Figura 1: Geometría del sistema bistático
\documentclass[border=5pt]{standalone}
\usepackage{tikz}
\usepackage{tikz-3dplot}
\usetikzlibrary{arrows.meta, decorations.pathreplacing, calc} % <-- ADICIONAR calc

\begin{document}
\tdplotsetmaincoords{70}{110}

\begin{tikzpicture}[tdplot_main_coords, scale=0.8]

% Definir puntos
\coordinate (Tx) at (2, 4, 3);
\coordinate (Rx) at (-3, -2, 3);
\coordinate (Target1) at (0, 0, -2);
\coordinate (Target2) at (1, 1, -2);
\coordinate (Surface1) at (0.8, 1.6, 0);
\coordinate (Surface2) at (-0.6, -0.4, 0);

% Plano del suelo
\fill[gray!20, opacity=0.5] (-4,-4,0) -- (4,-4,0) -- (4,4,0) -- (-4,4,0) -- cycle;
\draw[gray] (-4,-4,0) -- (4,-4,0) -- (4,4,0) -- (-4,4,0) -- cycle;

% Subsuelo
\fill[brown!30, opacity=0.3] (-4,-4,-3) -- (4,-4,-3) -- (4,4,-3) -- (-4,4,-3) -- 
                              (-4,-4,0) -- (4,-4,0) -- (4,4,0) -- (-4,4,0) -- cycle;

% Transmisor
\draw[blue, thick] (Tx) circle (0.15);
\fill[blue] (Tx) circle (0.1);
\node[above, blue] at (Tx) {Transmitter};

% Receptor
\draw[red, thick] (Rx) circle (0.15);
\fill[red] (Rx) circle (0.1);
\node[above, red] at (Rx) {Receiver};

% Targets
\foreach \target in {Target1, Target2} {
    \draw[black, thick] (\target) circle (0.1);
    \fill[black] (\target) circle (0.08);
}
\node[below] at (Target1) {Target 1};
\node[below] at (Target2) {Target 2};

% Trayectorias de rayos
% Tx → Surface1 → Target1
\draw[blue, thick, ->] (Tx) -- (Surface1);
\draw[blue, thick, dashed, ->] (Surface1) -- (Target1);

% Target1 → Surface2 → Rx
\draw[red, thick, dashed, ->] (Target1) -- (Surface2);
\draw[red, thick, ->] (Surface2) -- (Rx);

% Etiquetas de distancias
\node[midway, above] at ($(Tx)!0.5!(Surface1)$) {$R_1$};
\node[midway, right] at ($(Surface1)!0.5!(Target1)$) {$R_2$};
\node[midway, left] at ($(Target1)!0.5!(Surface2)$) {$R_3$};
\node[midway, below] at ($(Surface2)!0.5!(Rx)$) {$R_4$};

% Puntos de intersección
\fill[green] (Surface1) circle (0.05);
\fill[green] (Surface2) circle (0.05);

% Ejes de coordenadas
\draw[->] (0,0,0) -- (2,0,0) node[right] {$X$};
\draw[->] (0,0,0) -- (0,2,0) node[above] {$Y$};
\draw[->] (0,0,0) -- (0,0,2) node[above] {$Z$};

% Etiquetas de medios
\node[right] at (3, 3, 1.5) {Air ($n_1 = 1$)};
\node[right] at (3, 3, -1.5) {Ground ($n_2 > 1$)};

\end{tikzpicture}
\end{document}